\section{Statutory Implications}
\label{sec:statutory}

\subsection{One Verifiable Plan vs.\ 10,000 Unverifiable Plans}

Ensemble evidence is powerful for \emph{detection} but weak for \emph{selection}.
A court can use an ensemble to determine that a contested plan is an outlier.
But which plan should replace it?  The ensemble offers 10,000 candidates,
none of which has a principled claim over the others.  A special master tasked
with drawing a remedial map cannot simply ``pick the median ensemble plan''
without introducing a new layer of discretion.

The \ar{} plan solves this problem.  Given the census data and the seat count,
there is exactly one \ar{} plan per state.  It is:

\begin{itemize}
  \item \textbf{Verifiable}: Anyone with access to the public census data
    and the \texttt{redist} binary can reproduce it exactly, bit-for-bit.
  \item \textbf{Deterministic}: No human judgment enters after the inputs
    are fixed.
  \item \textbf{Politically inert}: No partisan registration data, no
    electoral history, no demographic targets beyond population equality.
  \item \textbf{Constitutionally grounded}: The \ar{} algorithm derives
    from Article~I~\S2 via the prime factorization of the Huntington-Hill
    seat count \citep{dellaLibera2026b11}.
\end{itemize}

An ensemble plan cannot make the first three claims.  Even if a court selected
the 50th-percentile plan on compactness, that plan reflects one random draw
and would change with a different seed, a different run, or a different
ensemble construction.

\subsection{The \texorpdfstring{\rucho{}}{Rucho} Framework and Algorithmic Plans}

\rucho{} held that partisan gerrymandering claims present a political question
not cognizable in federal court.  The majority distinguished between identifying
a gerrymander (possible) and providing a judicially manageable standard for
remediation (not possible under the Constitution).  The dissent argued that
ensemble evidence provides exactly such a standard.

The \ar{} plan offers a different response: bypass the manageability problem
entirely by providing a plan that requires no judicial discretion to select.
The court does not choose among plans; it certifies the \ar{} plan as
constitutionally required by the Huntington-Hill apportionment process.
This argument does not depend on \rucho{} being overruled; it operates at the
Article~I~\S2 level rather than the Equal Protection or First Amendment level.

\subsection{Ensemble Evidence as Corroboration}

Within the G-track framework, ensemble evidence serves as \emph{corroboration}
for the \ar{} plan rather than as its foundation.  The argument structure is:

\begin{enumerate}
  \item The \ar{} plan is constitutionally required (B.11, B.02).
  \item The \ar{} plan is not a partisan outlier: its partisan percentile
    is near the ensemble median (G.1, G.2).
  \item The \ar{} plan is compact but not an extreme outlier on compactness:
    its compactness percentile is $65$--$75$th (G.3).
  \item The \ar{} plan converges reproducibly: \cs{} at $T = 600$ certifies
    optimality (B.16).
\end{enumerate}

Claims 2 and 3 are corroborative: they show that the \ar{} plan possesses
properties independently associated with non-gerrymandered maps.  They do not
constitute the primary legal argument.  This ordering is important: ensemble
evidence is used to verify that the deterministic plan is reasonable, not to
select the plan.
